\chapter{Methodology}
\label{chap:methodology}

The methodology combines a Design Science Research-oriented artifact with
historical supervised learning. This approach fits the development and
evaluation of a software artifact for a practical engineering problem
\cite{hevner2004design,peffers2007design}. The artifact, OSS Risk Radar, supports
open-source dependency maintenance and software supply-chain risk triage
\cite{ossRiskRadarReadme}.

The training dataset uses real GitHub-hosted repositories and historical GitHub
events, not manually assigned or synthetic labels
(Section~\ref{sec:intro-scope}). For each repository, the builder creates
observation snapshots whose features were available at or before date \(t\).
The label comes from the following 12 months,
\((t, t + 12\text{ months}]\). The resulting task estimates future inactivity
from information available at the observation date.

\section{Research Design}
\label{sec:method-research-design}

The research design has three parts: OSS Risk Radar ingests dependency or
repository information and produces inactivity-risk scores from public signals;
a historical dataset supplies supervised examples; and quantitative evaluation
measures discrimination and calibration, complemented by formative manual case
review.

The artifact estimates inactivity risk from observable maintenance and activity.
It is neither a vulnerability scanner nor a definitive trust score: inactive
repositories may have no vulnerability, while active ones may still have
security weaknesses. The output prioritizes repositories for manual review.

OpenSSF Scorecard is a \emph{separate} security-posture signal in the enrichment
workflow \cite{zahan2023scorecard}. A Scorecard-derived quantity does \emph{not}
enter the observation-time feature set in
Section~\ref{sec:method-feature-engineering}.
Security-practice indicators appear as parallel analyst context, and the
predictive target remains restricted to activity and maintenance signals.

For each repository, OpenSSF provides an aggregate in \([0,10]\) and checks for
branch protection, code review, signed releases, least-privilege workflow tokens,
and known vulnerabilities \cite{zahan2023scorecard}. The service adjusts strong
and weak checks and returns a 0--100 \texttt{security\_posture\_score}. The
separate field is marked not model-derived when Scorecard is missing; it is not
imputed.\kiref{Claude Opus 5/1}

\section{Iterative Artifact Development}
\label{sec:method-iterative-development}

The first iteration validated snapshot construction, training, artifact export,
and manual scoring with a small pilot run that contributes no reported result.
Manual review found implausible scores: one operationally active project ranked
too high, while lower scores in other cases seemed to require more contextual
assessment. This formative error analysis led to separate features for activity,
releases, contributor dynamics and concentration, responsiveness, and backlog
pressure. It followed a build--evaluate--diagnose--refine cycle.

The second iteration produced the dataset and model artifacts used for the reported evaluation, expanding the empirical basis to the repository-first foundation dataset of Section~\ref{sec:method-data-sources}, whose seed buckets serve only as sampling provenance while the label is derived from future repository behavior.

\section{Definitions and Prediction Target}
\label{sec:method-definitions}

The unit of analysis is a GitHub repository. Each row represents a repository at
observation date \(t\). The binary prediction task is:

\begin{quote}
Given the observable repository and package state at time \(t\), predict whether the repository becomes inactive within the following 12 months.
\end{quote}

The observation date lies far enough in the past for the following twelve months
to have \emph{already elapsed}. A snapshot can only receive a supervised label after this
outcome window closes. Observation dates run quarterly from
\texttt{2023-01-01} to \texttt{2025-04-01}, and the GH Archive files extend to
\texttt{2026-06-16}. The final label window closes on
\texttt{2026-04-01}, within the recorded period
(Section~\ref{sec:method-observation-window}).

Each snapshot splits the historical archive at \(t\) into two disjoint periods:

\begin{itemize}
    \item \textbf{Backwards from \(t\)} lies the evidence the model is allowed to see. Observation-time features aggregate the event history over the 30, 90, and 365 days before \(t\) (Section~\ref{sec:method-feature-engineering}). This is the \emph{input}.
    \item \textbf{Forwards from \(t\)} lies the twelve-month window from which the label is computed, and which the model never sees for that row. This is the \emph{prediction target}.
\end{itemize}

Here, \emph{future} is relative to \(t\): the window follows the observation date
but precedes dataset construction. Its outcome was unknown at \(t\), although it
is available for supervised training later. Live scoring sets \(t\) to the
current date, for which the outcome window has not elapsed and no label
exists.\kiref{Claude Opus 5/4}

The exported \emph{inactivity label},
\texttt{label\_inactive\_12m}, is the negation of the future-window maintenance
variable \texttt{maintained\_12m}. The archive records positive evidence
\emph{of} maintenance: commits, contributors, releases, and merged pull requests. Defining
these four observable checks first makes the rule easier to verify and vary
(Section~\ref{sec:results-label-sensitivity}). Negation is exact, so the two
directions are equivalent. Only inactivity, the \emph{absence} of that evidence,
is exported and modeled as the positive class because the target is the risk
that maintenance stops.\kiref{Claude Opus 5/4}
Rows with incomplete future windows remain unlabeled and are excluded from
fitting.

The future label window is defined as:

\[
(t, t + 12\text{ months}]
\]

Figure~\ref{fig:observation-window} shows this temporal structure: the observation-time features summarize repository history up to \(t\) (Section~\ref{sec:method-feature-engineering}), and the label is derived only from the window after it.

\begin{figure}[ht]
\centering
\begin{tikzpicture}[>={Stealth[length=2.2mm]}]
  \fill[black!12] (-6,-0.32) rectangle (0,0.32);
  \fill[black!5]  (0,-0.32) rectangle (5,0.32);
  \draw[->] (-6.7,0) -- (7.4,0) node[right]{time};
  \draw[dashed,very thick] (0,-1.25) -- (0,1.7);
  \node[above] at (0,1.7){observation date $t$};
  \draw[dashed] (6,-0.9) -- (6,1.2);
  \node[above,align=center] at (6,1.2){\footnotesize archive\\[-2pt]\footnotesize coverage horizon};
  \node[align=center] at (-3,1.12){\footnotesize observation-time features};
  \node[align=center] at (-3,0.66){\footnotesize (trailing 30/90/365\,d,\ only data $\le t$)};
  \node[align=center] at (2.5,1.12){\footnotesize label window};
  \node at (2.5,0.66){\footnotesize $(t,\,t+12\text{ mo}]$};
  \draw (-6,0.32) -- (-6,-0.5) node[below]{\footnotesize $t-365$\,d};
  \draw (5,0.32) -- (5,-0.5) node[below]{\footnotesize $t+12$\,mo};
  \node[align=center,text width=5.2cm] at (2.5,-1.6){\footnotesize labeled only if the label window lies within the coverage horizon};
\end{tikzpicture}
\caption[Temporal structure of a repository observation snapshot]{Temporal structure of a repository observation snapshot. The whole diagram lies in the past: the observation date~\(t\) is chosen more than twelve months back, so that the label window after it has already elapsed and is recorded in the archive. Observation-time features are computed only from data at or before~\(t\), and the 12-month label only from the window after it; the dashed line at~\(t\) is the temporal leakage boundary (Section~\ref{sec:method-leakage-prevention}). A snapshot is labeled only when the archive-wide coverage horizon spans its full label window (Section~\ref{sec:method-observation-window}).}
\label{fig:observation-window}
\end{figure}

The current implementation treats a repository as maintained within the 12-month horizon if it is not archived or deleted by \(t + 12\text{ months}\) and if at least two of the following future-window activity checks are fulfilled:

\begin{enumerate}
    \item human commits occur in at least three distinct months,
    \item at least two contributors are active,
    \item at least one release or package version publication occurs,
    \item at least two pull requests are merged.
\end{enumerate}

This is the \emph{two-of-four maintenance rule}. Requiring two checks prevents a
single drive-by commit or automated release from establishing maintenance.
Requiring all four would misclassify healthy projects that do not release or
receive external pull requests during a given year. The rule remains measurable
from public history.

Merged pull requests proxy broader maintainer interaction. GH Archive records
typed events, not resolved roles \cite{gharchive}; identifying maintainer
comments would require attribution logic and remains vulnerable to bot
distortion \cite{hasan2023firstresponse}. Merge events are unambiguous in the
event stream. The artifact documents the implementation
\cite{ossRiskRadarDatasetBuilder}.

The offline build supplies no \texttt{archived\_at} timestamp
(Section~\ref{sec:method-offline-metadata}), so the archival condition never
fires and the label uses only the two-of-four rule. Archived repositories
normally receive none of the four activity signals and are already marked
inactive, making this a conservative simplification. The condition would,
however, miss a repository archived while still receiving activity.

The two-of-four requirement and the cut-offs for months, contributors, releases,
and merged pull requests are heuristic design choices, not externally validated
constants. They require sustained activity. Because they affect class
prevalence, Section~\ref{sec:results-label-sensitivity}
(Appendix~\ref{appendix:label-sensitivity}) repeats the evaluation with stricter
and more permissive variants.

\section{Data Sources}
\label{sec:method-data-sources}

GH Archive is the primary historical source. It distributes public GitHub event
streams as hourly archives \cite{gharchive}, which the builder parses into
normalized histories of commits, issues, pull requests, releases, stars, and
forks.

deps.dev supplies package and dependency metadata when package-to-repository
mapping is required \cite{depsdevApi}. GitHub supplies stable fields such as
creation time, default branch, and fork state. OpenSSF Scorecard remains a
separate security-practice source \cite{zahan2023scorecard}
(Section~\ref{sec:method-research-design}). Present-day GitHub activity totals
are excluded from historical snapshots because they would leak information from
after the observation date.

The first construction stage produces the \emph{foundation seed}, a fixed list
of 5{,}000 repository identities without features or labels. The second stage
replays each history at ten observation dates. Storing the seed makes reruns use
the same repositories instead of a new search result.\kiref{Claude Opus 5/4}

The seed contains public, non-fork GitHub repositories with a license object and
at least 100 stars by default. Separate active, dormant, and archived buckets
control the mixture. They are neither features nor labels: labels come
exclusively from activity in the 12-month future window
\cite{ossRiskRadarDatasetBuilder}. They are, however, defined on repository state
after the observation dates, which is qualified below. The target mixture is
45\% active and 55\% dormant or archived.

The dormant and low-star active buckets could not be filled, so a fallback added
active repositories. The realized seed is 52.0\% active, 28.0\% dormant, and
20.0\% archived (Appendix~\ref{sec:appendix-seed-composition}), so the intended
oversampling is only partly realized at seed level. The held-out inactive rate
reported in Chapter~\ref{chap:evaluation} is therefore more relevant than the
seed buckets. It exceeds
reported unsampled rates: Coelho et al.\ observe that \(16\%\) of \(2{,}927\)
active projects became unmaintained within one year, while Miller et al.\ report
\(15\%\) abandonment among widely used npm packages over six years
\cite{coelho2020maintained,miller2025abandonment}. The probabilities are
calibrated to the sampled prevalence, not a deployment population
(Sections~\ref{sec:method-evaluation-strategy} and~\ref{sec:method-limitations}).

\label{sec:method-retrospective-frame}
\textbf{The frame is sampled retrospectively.} The seed was generated on
\texttt{2026-06-11}, after every label window had closed, and its bucket queries
read repository state at that date. The active buckets require a push within the
preceding year (\texttt{pushed:>=2025-06-11}), the dormant buckets require no
push for two years (\texttt{pushed:<=2024-06-11}), and the archived buckets
require the archive flag to be set. Entry into the frame is therefore decided by
behavior that falls partly inside the outcome windows the model predicts. The
buckets do not enter the feature vector and do not define the label, but they are
not independent of it either. In the final test cohort, the inactive rate is
\(1.000\) across the \(1{,}399\) dormant rows and \(0.942\) across the
\(1{,}000\) archived rows. It is \(0.409\) across the \(2{,}095\) active rows
and \(0.209\) across the \(506\) fallback rows. The dormant rate reaches
\(1.000\) from the \texttt{2024-07-01} cohort onward.

Two consequences follow, and both restrict how the reported numbers may be read.
The inactive rate of \(0.661\) describes this construction rather than a base
rate obtainable by sampling repositories at \(t\). The aggregate task is also
easier than the deployment task. Much of it consists of
separating buckets that were formed from later behavior and are already legible
at \(t\) through commit recency. Consistently, the trivial recency baseline of
Section~\ref{sec:results-baselines} reproduces most aggregate ranking skill. All
reported metrics are conditional on this retrospectively
sampled frame. Chapter~\ref{chap:validity} returns to the threat, and drawing the
frame from repository state at each observation date, rather than at generation
time, is recorded as future work in
Chapter~\ref{chap:conclusion}.\kiref{Claude Opus 5/11}

\label{sec:method-offline-metadata}
\textbf{Offline repository metadata and its consequences.} The reported build
uses the stored seed and GH Archive files without live GitHub, registry, or
deps.dev calls. It is repeatable without credentials, rate limits, or changes in
mutable metadata. All observation-time features are reconstructed from event
history, so the predictive content is unaffected. The offline path leaves
creation date, archival timestamp, and fork flag unpopulated. This makes
\texttt{repo\_age\_days} identically zero and the archival label condition and
fork filter inert. Appendix~\ref{sec:appendix-undefined-features}
details these effects; Chapter~\ref{chap:conclusion} lists metadata enrichment as
future work.

\section{Dataset Construction Workflow}
\label{sec:method-dataset-construction}

Figure~\ref{fig:dataset-pipeline} shows the construction sequence. The builder
loads the foundation seed, resolves repository identities, enriches stable
metadata, normalizes GH Archive events, and creates quarterly snapshots. For each
snapshot, it computes features from data at or before \(t\) and the 12-month
label from \((t, t + 12\text{ months}]\). It exports both the training table and
a feature cache for runtime full-history scoring.

Feature and label computation read disjoint parts of the normalized timeline and
meet only in the exported row. Appendix~\ref{appendix:dataset-construction} lists
the intermediate files that make each stage inspectable. The final table contains
one snapshot per row, the features from
Section~\ref{sec:method-feature-engineering}, and one label column. It matches the
existing training interface, so reconstructed data use the same training path as
other runs. Appendix~\ref{sec:method-reproducibility} gives the rebuild commands,
run documentation, and dataset hash.

\begin{figure}[htbp]
\centering
\begin{tikzpicture}[
  font=\footnotesize,
  box/.style={draw,black!55,rounded corners=1pt,align=center,inner sep=3.4pt,minimum height=7.6mm,text width=21mm},
  src/.style={box,fill=black!8},
  feat/.style={box,fill=black!14},
  lab/.style={box,fill=black!5},
  sink/.style={box,fill=black!20},
  ar/.style={->,>=Stealth,black!65,line width=0.45pt}
]
  % --- ingestion chain -------------------------------------------------
  \node[src] (seed)  at (0,0)    {foundation seed\\(5{,}000 repos)};
  \node[src] (arch)  at (3.4,0)  {GH Archive\\event files};
  \node[src] (hist)  at (6.8,0)  {normalized\\event history};
  \node[src] (snap)  at (10.2,0) {quarterly\\snapshots at $t$};

  \draw[ar] (seed) -- (arch);
  \draw[ar] (arch) -- (hist);
  \draw[ar] (hist) -- (snap);

  % --- the split -------------------------------------------------------
  \node[feat] (feats) at (10.2,-2.0) {observation-time\\features};
  \node[lab]  (labs)  at (10.2, 2.0) {12-month\\label};

  \draw[ar] (snap) -- node[right,font=\scriptsize,pos=0.42,align=left]
    {\ events $\le t$} (feats);
  \draw[ar] (snap) -- node[right,font=\scriptsize,pos=0.42,align=left]
    {\ events $>t$} (labs);

  % leakage boundary
  \draw[dashed,black!70,line width=0.7pt] (7.9,-3.0) -- (7.9,3.0);
  \node[black!70,font=\scriptsize\itshape,align=center,anchor=south]
    at (7.9,3.0) {temporal leakage boundary};

  % --- outputs ---------------------------------------------------------
  \node[sink] (train) at (13.6,-1.0) {training snapshot\\table (50{,}000 rows)};
  \node[sink] (cache) at (13.6, 1.2) {repository\\feature cache};

  \draw[ar] (feats) -- (train);
  \draw[ar] (labs)  -- (train);
  \draw[ar] (feats) -- (cache);

  \node[font=\scriptsize,black!60,anchor=north,align=center]
    at (13.6,-1.85) {to model training\\(Section~\ref{sec:method-model-training})};
  \node[font=\scriptsize,black!60,anchor=south,align=center]
    at (13.6,1.95) {to runtime scoring\\(Chapter~\ref{chap:implementation})};
\end{tikzpicture}
\caption[Dataset construction pipeline and its leakage boundary]{Dataset
construction pipeline. The upper chain is ingestion: the fixed foundation seed
selects the repositories, GH Archive supplies their raw events, and these are
normalized into a per-repository chronological history from which quarterly
observation snapshots are cut. At each snapshot the pipeline splits. Everything at
or before the observation date \(t\) may become a \emph{feature}; everything after
it may only become the \emph{label}; nothing crosses the dashed boundary
(Section~\ref{sec:method-leakage-prevention}). The two paths meet only in the
exported row. The feature path is written a second time, without labels, as the
repository feature cache that lets the runtime score a previously seen repository
in the full-history regime. Appendix~\ref{appendix:dataset-construction} lists the
individual steps and the intermediate file written at each of them.}
\label{fig:dataset-pipeline}
\end{figure}

\section{Observation Window and Historical Coverage}
\label{sec:method-observation-window}

The reported run contains 50,000 rows: 5,000 repositories at ten quarterly dates
from \texttt{2023-01-01} to \texttt{2025-04-01}. Its GH Archive files span
\texttt{2021-01} to \texttt{2026-06-16}. The latest label window ends on
\texttt{2026-04-01}, so all 50,000 rows are labeled.

The time-ordered split assigns the oldest 37,500 rows to training, the next 7,500
to validation, and the newest 5,000 to testing. Because the ten cohorts have
equal size, the \(10\%\) test partition is the final observation date; every test
row is observed on \texttt{2025-04-01}.

The cut is made by row position rather than by cohort, so the boundary between
training and validation falls inside the \texttt{2024-10-01} cohort: \(2{,}500\)
of its rows are assigned to training and \(2{,}500\) to validation. Training thus
covers \texttt{2023-01-01} to \texttt{2024-10-01} and validation
\texttt{2024-10-01} to \texttt{2025-01-01}. Rows are ordered only by observation
date and ties keep their original order, so the division within that cohort is by
position and not by any repository property. The test partition is unaffected and
remains one complete cohort.

\label{sec:method-label-availability}
\textbf{Labels are used before they could have been observed.} Each label
requires the twelve months following its observation date. The latest training
rows are observed on \texttt{2024-10-01} and their labels close on
\texttt{2025-10-01}; the latest validation labels close on \texttt{2026-01-01}.
Both lie after the test observation date of \texttt{2025-04-01}. An analyst
standing at \texttt{2025-04-01} could therefore not have fitted this model,
because part of its training signal had not yet occurred. The evaluation is a
retrospective chronological one with overlapping outcome windows rather than a
simulation of prospective deployment. It supports the claim that the model ranks
later observations correctly when fitted on earlier ones; it does not establish
that the same model could have been built in time to issue those predictions.
Removing the overlap would require an embargo of one full label horizon, which
here means discarding the three cohorts whose outcome windows extend past the
test date --- \texttt{2024-07-01}, \texttt{2024-10-01}, and \texttt{2025-01-01},
or \(15{,}000\) of the \(45{,}000\) non-test rows. With ten cohorts available
this was judged too costly for the reported run and is recorded as future work
(Chapter~\ref{chap:conclusion}).\kiref{Claude Opus 5/11}

The model is fitted on earlier observations and tested on later ones, matching
the forecasting setting. However, each repository contributes ten snapshots and
therefore appears in training and testing at different dates. Performance could
partly reflect recognition of seen repositories. A repository-disjoint
evaluation assigns every snapshot of a repository to one partition, leaving
only unseen repositories in its test set. Appendix~\ref{sec:impl-grouped-split}
describes this split, and Section~\ref{sec:results-robustness} reports the result.

A snapshot is labeled only if the archive covers the full twelve months \emph{after}
\(t\), shown on the right of Figure~\ref{fig:observation-window}. Other rows are
retained but excluded from fitting. Because inactivity is an absence of events,
incomplete coverage could otherwise turn missing data into positive labels.

This run contains no such rows: coverage ends on \texttt{2026-06-16}, after the
latest label window closes on \texttt{2026-04-01}. An observation after roughly
\texttt{2025-06}, or missing archive files, would activate the gate. It removes
no rows from the reported dataset.\kiref{Claude Opus 5/4} A smaller subset is
used only for local pipeline smoke tests and never for reported results.
Thus, all \(50{,}000\) reported rows pass this coverage gate.

\section{Feature Engineering}
\label{sec:method-feature-engineering}

Feature set \texttt{feature\sdash set\sdash v4\sdash full\sdash history} contains
60 inputs derived before \(t\). Of these, 43 reconstruct commit and contributor
activity, issue and pull-request flow, releases, risk proxies, responsiveness,
and backlog pressure. The other 17 describe the current snapshot, including
dependency context, popularity, and ecosystem. The cold-start regime uses these
17 features (Section~\ref{sec:method-prediction-regimes}).

\emph{Backlog pressure} comprises three historical quantities: relative
open-issue growth over the preceding 90 days, issues older than 90 days that
remain open at \(t\), and their share of the open backlog. The cold-start set also
contains backlog size and 90-day open-issue growth. Together they indicate
whether open items are growing and ageing independently of response speed.

Because the seed is repository-first and the offline build makes no registry
call, ecosystem indicators and dependency-context flags are constant. Eight of
the 60 inputs are constant across all 50,000 rows and are inert, because
  zero-variance columns cannot split a tree. Two further inputs,
  \texttt{package\sdash age\sdash days} and
  \texttt{versions\sdash published\sdash 365d}, vary only with the observation
  date and not between repositories, so ten inputs in total carry no
  repository-level signal and roughly 50 remain informative. The two counts describe the
  implemented vectors, not the number of informative signals. Appendix~\ref{sec:appendix-undefined-features}
lists the affected features and explains why the two cohort-dependent inputs
cannot change the reported ROC-AUC.\kiref{Claude Opus 5/6}
Appendix~\ref{appendix:health-indicators}
(Table~\ref{tab:background-health-indicators}) maps the health-indicator
categories of Section~\ref{sec:background-oss-health} to feature groups.
Appendix~\ref{appendix:feature-reference}
(Tables~\ref{tab:method-feature-groups}, \ref{tab:method-feature-reference},
and~\ref{tab:method-snapshot-features}) defines every feature.

The builder parses GH Archive events into chronological histories of commits,
issues, pull requests, releases, stars, and forks \cite{gharchive}. Features are
counts over the 30, 90, and 365 days before \(t\), state at \(t\), or derived
ratios, latencies, and drop-off measures. They lie before the leakage boundary in
Figure~\ref{fig:observation-window}. Trailing windows capture trajectory instead
of lifetime volume.

Commit and contributor counts apply a \emph{human-actor filter} that excludes
events attributed to bot accounts. It prevents automated activity, such as
weekly dependency-bot commits, from appearing as maintenance. The label uses the
same filter, and
Appendix~\ref{appendix:bot-ablation} measures its effect by rebuilding the
dataset without it. Relative and trend-based features supplement absolute counts
that do not compare across project sizes.

Revision activity predicts project survival \cite{robinson2022survival}, which
motivates the commit, contributor, and release windows. Issue, pull-request, and
responsiveness signals follow the CHAOSS project-health vocabulary
\cite{chaossStarterProjectHealth,goggins2021transparent}. Since individual
context-free indicators are insufficient across heterogeneous projects
\cite{yehudi2023contextfree}, the model learns a calibrated combination.

The responsiveness features are \emph{reconstructed proxies}, not the canonical
CHAOSS \emph{Time to First Response} metric. They follow thesis-specific rules
over GH Archive events. A first response is the earliest comment from an account
other than the person who opened the issue or pull request. This account need not
be a maintainer and may be a bot
\cite{hasan2023firstresponse}. For pull requests without comments, the merge or
close event substitutes for a response.

Reconstructed star and fork counts use a second sense of \emph{proxy}: they are
lower bounds when
archive coverage begins after repository creation. These limitations do not
invalidate the features, but attributions describe observable behavior rather
than canonical CHAOSS measurements.\kiref{Claude Opus 5/4}

Reconstructed issue backlogs likewise describe the archived tracker, not a full
API result.

\label{sec:method-undefined-features}
\textbf{Undefined versus zero-valued features.} Zero is the \emph{most favorable} value
for several variables and must not encode missing data. When no commit is
observed, ``days since last commit'' is \emph{right-censored} at the age of the
observable record. This represents maximal observed staleness instead of a
recent commit. A metric that does not exist, such as median merge time for a
repository with no merged pull request, remains \emph{undefined} for model-specific
imputation. Appendix~\ref{sec:appendix-undefined-features} defines the censoring
bound, imputation, and two ``no events'' indicators.

This distinction keeps data availability from masquerading as maintenance
behavior in the attributions of Section~\ref{sec:results-attribution}. A
zero-valued latency would make silence look like immediate responsiveness. The
revised feature set also excludes unreliable point-in-time predictors
(Section~\ref{sec:method-iterative-development}). For example,
\texttt{repo\_archived\_at\_obs} remains diagnostic metadata while
already-archived rows are excluded from fitting, preventing archival from
becoming a shortcut label.

\section{Full-History and Cold-Start Prediction Regimes}
\label{sec:method-prediction-regimes}

Each model is trained in two regimes. \emph{full-history} uses all 60 inputs,
including 43 reconstructed commit, release, contributor, issue, pull-request, and
activity-drop features before \(t\). \emph{cold-start} uses the 17
current-snapshot features and requires no event history.

This ablation on \emph{feature availability} measures the value of reconstructing
history from tens of thousands of public archive files compared with one metadata
read. Historical features are theoretically available for any public repository:
the regimes do not contrast what a scorer \emph{can} know with what it
\emph{cannot}.
However, \(7{,}687\) of \(50{,}000\) rows have no earlier observed event, leaving
their historical features censored or undefined
(Section~\ref{sec:method-undefined-features}). At runtime, repositories absent
from the feature cache also use cold-start, although this is not the motivation
for the ablation. The comparison asks what reconstruction is \emph{worth}.

Across the full test set, full-history exceeds cold-start by only \(0.005\)
ROC-AUC. Many repositories are already quiet at \(t\), so recency makes their
outcome easy to rank without reconstructed history. A repository with no commits
for two years, for example, needs no reconstructed trajectory to appear at risk.
Section~\ref{sec:results-history-ablation} shows that this depends on
\emph{which repositories the test set contains}: the difference is larger among
repositories that remain active at \(t\) and later decline.

\section{Leakage Prevention}
\label{sec:method-leakage-prevention}

Figure~\ref{fig:observation-window} shows the boundary at \(t\) between feature
construction and label derivation. The builder enforces:

\begin{enumerate}
    \item Observation-time features are derived only from timestamps \( \leq t \).
    \item Labels are derived only from timestamps in \( (t, t + 12\text{ months}] \).
    \item Rows with incomplete future-window coverage remain unlabeled.
    \item Live GitHub enrichment is limited to stable repository metadata such as creation time, default branch, and fork state.
    \item Historical stars, forks, issues, pull requests, and releases are not taken from current GitHub API totals.
    \item Seed buckets such as active, dormant, and archived are used only for sampling provenance and not as labels. They are drawn from repository state at seed-generation time and therefore constrain the frame rather than the row (Section~\ref{sec:method-retrospective-frame}).
\end{enumerate}

At \(t\), later activity is unknown. It is excluded from features and used only
to derive the label.

Activity persistence is distinct from leakage. Days since the last commit,
activity drop, and days since the last release use valid information available at
\(t\), but already-quiet repositories tend to remain quiet. Some discrimination
therefore reflects persistence instead of early decline. This motivates the
recency baseline and active-at-\(t\) evaluation in
Section~\ref{sec:method-evaluation-strategy}.

\section{Model Training}
\label{sec:method-model-training}

One notebook performs all fitting and exports the model files. The scoring
service only loads these files \cite{ossRiskRadarReadme}.
Section~\ref{sec:impl-deployment} describes this reproducibility boundary.

The learned families are Logistic Regression, XGBoost, and a feedforward neural
network, each trained in both regimes. Logistic Regression provides an
interpretable linear model. XGBoost supplies non-linear gradient-boosted trees
\cite{chen2016xgboost}; the neural network learns non-linear interactions through
hidden layers \cite{james2021islr}.

Two non-learned baselines rank repositories by days since the last human commit
or by negated annual commit volume. They select thresholds on validation data and
are evaluated on the test set. Commit volume covers all test rows, while recency
is defined for \(4{,}556\) rows (Section~\ref{sec:results-baselines}). Both run
inside the training notebook, linking their results to the same run as the
learned models. Because the label is activity-related, these baselines measure
how much ranking skill one activity signal already provides
(Section~\ref{sec:method-leakage-prevention}). Section~\ref{sec:results-baselines}
shows that the learned advantage is concentrated among repositories still active
at \(t\).

The exact model configurations and hyperparameters are documented in Appendix~\ref{sec:impl-from-scratch}; all learned models are calibrated on the validation set before evaluation.

Training uses a time-aware split with the default ratio:

\[
75\% \text{ training}, \quad 15\% \text{ validation}, \quad 10\% \text{ test}
\]

Rows are ordered by observation date. The earliest rows train the model;
validation predictions fit the histogram calibrator and select the decision
threshold. The final test set is reserved for reporting. This avoids mixing later
observations into earlier model development.

\section{Evaluation Strategy}
\label{sec:method-evaluation-strategy}

The primary evaluation uses the time-aware test set for neither fitting nor
calibration. Appendix~\ref{sec:method-metric-definitions} defines accuracy,
precision, recall, F1, Brier score, log loss, ROC-AUC, average precision,
expected calibration
error, and the combined quality score.

The inactive base rate is approximately 0.66, so ROC-AUC is a reasonable primary
ranking metric. Precision-recall summaries are more informative under strong class
imbalance \cite{saito2015precision}, which is why average precision is reported
for the active-at-\(t\) subset defined below, where the inactive rate falls to
\(0.225\) (Section~\ref{sec:results-precision-recall}). Precision and recall are reported at the
selected operating point. Since precision depends on prevalence, it is
conditional on this sample and is not a deployment guarantee or headline
comparison. Calibration is assessed separately because probability scores must
match observed frequencies \cite{niculescu2005probabilities}.

The models are compared with recency and evaluated separately on repositories
still active at \(t\). This subset limits domination by activity persistence
(Section~\ref{sec:method-leakage-prevention}) and represents the harder case in
which early warning can still affect a decision. Active-at-\(t\) means at least one
bot-filtered human commit in the trailing 90 days:
\(\text{active-at-}t = \mathbf{1}\!\left(\text{\texttt{commits\_90d}} > 0\right)\).
The 90-day window is an operational choice, not a tuned threshold. A multi-signal
variant with recent releases or merged pull requests remains future work.

Because inactivity as defined in Section~\ref{sec:method-definitions} is an
operational proxy, formative manual review identifies
context-sensitive cases and predictions made for the wrong reason; it drove the
first feature revision. A systematic case review is future work
(Chapter~\ref{chap:conclusion}). The held-out split is an \emph{internal}
evaluation of temporal generalization. Chapter~\ref{chap:validity} discusses
external validity.

The staged full-history artifact is
\texttt{xgboost\sdash full\sdash history}, version \texttt{0.2.0}, trained with
\texttt{feature\sdash set\sdash v4\sdash full\sdash history}.
Chapter~\ref{chap:evaluation} compares its held-out metrics with Logistic
Regression, \texttt{neural\sdash net\sdash full\sdash history}, and the two
baselines.

\section{Methodological Limitations}
\label{sec:method-limitations}

This section records data and design limitations; Chapter~\ref{chap:validity}
discusses the broader validity threats.

Missing GH Archive hours leave histories incomplete. Public GitHub events omit
development elsewhere, off-platform releases, and stable software requiring
little visible activity. Package-to-repository mappings can be ambiguous after
moves or with incomplete metadata. The inactivity label is an operational proxy,
and its heuristic thresholds affect prevalence
(Section~\ref{sec:method-definitions}).

The seed covers visible, licensed repositories above a minimum star count, not
all of GitHub. Oversampling dormant and archived repositories makes
probabilities conditional on the sampled prevalence; deployment risk estimates
require recalibration (Section~\ref{sec:method-data-sources}).

The internal split tests temporal generalization within the dataset. External
validation on a newly sampled repository set remains future work
(Chapter~\ref{chap:conclusion}).
